% sup4_07_S21b: Supplement part 4, step S21, II: the one-loop classes (3.6)-(3.30) and the subtraction identities
% (3.19), (3.31), (3.34). Input through INPUTS_part4.tex.

\subsection{S21, II: the one-loop classes and the subtraction identities}\label{sup4:ssec:S21b}

Notation: $j_0$ is the lowest index of the external legs of the class, $\xi=L^{-j_0}$, and after \cref{sup4:lem:S214,sup4:lem:S215} every internal line is $\breve G_l=C^\xi+R_l$ with $R_l$ as in \cref{sup4:lem:R}; $\hat G_s$, $\hat G_{\mathrm{ph}}$ denote the rescaled regrouped scalar and photon lines. Sums over $y'$ carry the weight $\xi^3$. The bounds $|C^\xi|,|\hat G|\le C(\xi+\rho)^{-1}e^{-\delta\rho}$ (\cref{sup4:lem:C}, \cref{sup4:cor:Z1}), $|\partial R|\le C(1+|\log(\xi+\rho)|)e^{-\delta\rho}$ and $|\partial R\partial^*|\le C(\xi+\rho)^{-1}e^{-\delta\rho}$ (\cref{sup4:lem:R}) are used throughout.

\begin{lemma}[{scalar self-energies with $D=-1$, \cite[(3.6)--(3.17)]{Ba3}}]\label{sup4:lem:S218}
\textup{(a)} The two tadpoles of \cite[(3.6)]{Ba3} and the last two counterterms of \cite[(3.7)]{Ba3} cancel exactly (``Wick ordering''): the class graph has the loop line $G_k(0)(x,x)$ or the torus photon line at $(x,x)$ and the localisation of \cite[(1.6)]{Ba3} to unit cubes of $\Omega_1$, resp.\ the factors $g_k(x)^2$; the counterterms \cite[(1.29)]{Ba3} carry the same localisation and the same lines once the torus scalar line is identified with $G_k(0)$, the difference being an image term, an external field of norm $O(e^{-\delta S_k})$ (\cref{sup4:lem:F2}). \textup{(b)} For \cite[(3.8)]{Ba3} minus its counterterm, i.e.\ the expression \cite[(3.9)]{Ba3}: Taylor's formula (\cref{sup4:lem:S213}) applied to the leg with the smaller index \cite[p.~435]{Ba3} gives \cite[(3.11)]{Ba3}. Its second term is the generalised graph \cite[(3.12)]{Ba3}, of degree $-1+1+\alpha=\alpha$; the factor $|x'-x|^{1+\alpha}$ multiplies every subgraph containing a line of $G_0$, since both lines join $x$ and $x'$, so its \cite[(2.4)]{Ba3}-subgraph has degree $1+\alpha>0$. Its first term, regrouped (\cref{sup4:lem:S214}) and rescaled (\cref{sup4:lem:S215}, $D=0$), is the vertex \cite[(3.15)]{Ba3} with coefficient
\[
c_\mu(y)=\sum_{y'}\xi^3\sum_\nu q(\partial_\nu\hat G_s\partial_\nu^*)(y,y')\,q\,g(y)\hat G_{\mathrm{ph}}(y,y')g'(y')\,(y'_\mu-y_\mu),
\]
and $|c_\mu|\le C$ uniformly in $\xi$, $j_0$, $k$. This is \cite[(3.17)]{Ba3}.
\end{lemma}

\begin{proof}
(a) is the identity of the lines and localisations stated. (b): write $\hat G_s=C^\xi+R_s$. Since $C^\xi$ is translation invariant, $\sum_\nu\partial_\nu C^\xi\partial_\nu^*=-\Delta^\xi C^\xi=\delta^\xi-C^\xi$ and $\delta^\xi(y'-y)(y'_\mu-y_\mu)=0$, so
\[
c_\mu=-q^2\sum_{y'}\xi^3C^\xi(y-y')\,g\hat G_{\mathrm{ph}}g'\,(y'-y)_\mu+q^2\sum_{y'}\xi^3\Big(\sum_\nu\partial_\nu R_s\partial_\nu^*\Big)(y,y')\,g\hat G_{\mathrm{ph}}g'\,(y'-y)_\mu,
\]
which is \cite[(3.16)]{Ba3} (whose first line prints $g(y')$ for $g'(y')$). The first sum is $\le C\sum\xi^3(\xi+\rho)^{-2}\rho e^{-\delta\rho}\le C$, and by \cref{sup4:lem:R} the second is bounded in the same way.
\end{proof}

\begin{lemma}[{scalar self-energies with $D=0$, \cite[(3.18)--(3.20)]{Ba3}}]\label{sup4:lem:S219}
If $\Sigma(x,x')$ is the expression of a graph $G_0$ of the class with external legs at $x$, $x'$, and its counterterm graph is $G_0$ with both legs at $x$ (\cite[p.~417]{Ba3}: ``represented by the same graph $G$ but with both external legs localized in $x$''; so it has the same lines, slices and orderings), then \cite[(3.19)]{Ba3} holds exactly:
\[
\sum\eta^{2d}\phi(x)\cdot\Sigma(x,x')\phi'(x')-\sum\eta^d\phi(x)\cdot\Big(\sum_{x'}\eta^d\Sigma(x,x')\Big)\phi'(x)=\sum\eta^{2d}\phi(x)\cdot\big(\Sigma(x,x')|x'-x|^\alpha\big)\frac{\phi'(x')-\phi'(x)}{|x'-x|^\alpha}.
\]
Along a chain of vertices $x=x_0,\dots,x_l=x'$ of $G_0$, $|x'-x|^\alpha\le\sum_i|x_{i+1}-x_i|^\alpha$ (concavity), which distributes the gain $+\alpha$ to the lines of a path from $x$ to $x'$. A subgraph containing no line of the path is a proper subgraph of $G_0$, hence positive when $G_0$ is primitively divergent within the ordering; the H\"older quotient of the external leg costs $\alpha$ on that leg, and $D\ge\frac12>\alpha$ for every positive subgraph containing it.
\end{lemma}

\begin{proof}
The identity is the definition of the counterterm graph; the degree bookkeeping is as stated.
\end{proof}

\begin{lemma}[{the differentiated-leg graphs \cite[(3.21)--(3.24)]{Ba3}}]\label{sup4:lem:S2110}
\textup{(a)} The first graph of \cite[(3.21)]{Ba3} (a \cite[(1.8)]{Ba3}-vertex with $n=2$ whose two $A'$-legs are contracted) has a coefficient of the type $\eta e_k^2G_{\mathrm{ph}}(b,b)$, bounded uniformly: by \eqref{sup4:eq:Z1} at $\rho=0$, $|G(x,x)|\le C_1\eta^{-1}$ in $d=3$; this is the printed ``$\eta G_k(x,x)$ is convergent''. \textup{(b)} The second graph, after Taylor at order $0$, is \cite[(3.22)]{Ba3}: a remainder of degree $\alpha$ plus the vertex \cite[(3.23)]{Ba3}, whose coefficient (of the vertex $\phi\cdot\partial_\mu\phi'$) is bounded uniformly in $\xi$, $j_0$, $k$.
\end{lemma}

\begin{proof}
(b) Write $g'(x')=g'(x)+(g'(x')-g'(x))$; the difference carries $|\partial g'||x'-x|$ (degree $+1$) and is convergent. The main term, rescaled and with $\hat G=C^\xi+R$, is $q^2g(x)g'(x)[\sum_y\xi^3(\partial_\mu C^\xi)(y)C^\xi(y)+(\text{terms with an }R)]$. Exactly, $\sum_y\xi^3(\partial_\mu C^\xi)(y)C^\xi(y)=-\frac\xi2\sum_y\xi^3|(\partial_\mu C^\xi)(y)|^2$, by symmetrising with $\sum_yC(y+\xi e)^2=\sum_yC(y)^2$; by \cref{sup4:lem:C} this is $\le\frac\xi2C\sum\xi^3(\xi+|y|)^{-4}\le C$. The $R$-terms are $\le C\sum\xi^3[(1+|\log(\xi+\rho)|)(\xi+\rho)^{-1}+(\xi+\rho)^{-2}]e^{-\delta\rho}\le C$ by \cref{sup4:lem:R,sup4:lem:C}. (\cite[(3.24)]{Ba3} shows a finite limit; only boundedness is used.)
\end{proof}

\begin{lemma}[{vector self-energies with $D=-1$, \cite[(3.25)--(3.30)]{Ba3}}]\label{sup4:lem:S2111}
Let the class be all order-$e_k^2$ graphs with two external vector legs: the two bubbles and the two tadpoles of \cite[(3.25)]{Ba3}. Apply Taylor to first order on the leg $g'A'$ at $x'$ around $x$ (\cite[(3.26)]{Ba3}). The H\"older remainder is a generalised graph of degree $-1+1+\alpha=\alpha$. After regrouping and rescaling the zero-order coefficient is $\Pi^{(\eta)}_{\mu\mu'}(x)=(L^{j_0}\eta)^{-1}\Pi^{(\xi)}_{\mu\mu'}(y)$ and the first-moment coefficient is $\Pi^{(\eta)}_{\mu\mu'\nu}(x)=\Pi^{(\xi)}_{\mu\mu'\nu}(y)$, with lines $\breve G_l=C^\xi+R_l$. Then: \textup{(i)} the all-$C$ part of $\Pi^{(\xi)}_{\mu\mu'}$ vanishes (\cref{sup4:lem:S216}(d); Ba{\l}aban's \cite[(3.29)]{Ba3} via \cite[(2.26)]{Ba3} at $M^2=1$), and the all-$C$ part of $\Pi^{(\xi)}_{\mu\mu'\nu}$ vanishes (\cref{sup4:lem:S217}(b); \cite[(3.27)--(3.28)]{Ba3}); \textup{(ii)} the terms with at least one $R$-line are bounded uniformly in $\xi$, $j_0$, $k$, both for $\Pi^{(\xi)}_{\mu\mu'}$ and for $\Pi^{(\xi)}_{\mu\mu'\nu}$. So \cite[(3.30)]{Ba3} holds: the class equals two generalised graphs of degree $\alpha$, plus the vertex $A\cdot\Pi_{\mu\mu'\nu}\partial A$ with a bounded coefficient, plus the vertex $A\cdot\Pi_{\mu\mu'}A$ with coefficient $(L^{j_0}\eta)^{-1}O(1)$. The last is a zero-momentum remainder vertex of power $-1$ at the scale $L^{j_0}\eta$ of the external legs; it is kept and bounded, not claimed to vanish.
\end{lemma}

\begin{proof}
(i) is as stated. (ii) The $R$-terms of the bubbles are of four kinds. $(C\partial^*)(R\partial^*)$ and $(R\partial^*)(R\partial^*)$: bounded by $C\sum\xi^3(\xi+\rho)^{-2}(1+|\log(\xi+\rho)|)e^{-\delta\rho}\le C$. $C\cdot(\partial R\partial^*)$: bounded by $C\sum\xi^3(\xi+\rho)^{-2}e^{-\delta\rho}\le C$. $R\cdot(\partial R\partial^*)$: bounded likewise. $R(y,y')(\partial_{\mu'}C\partial_\mu^*)(y'-y)$: here absolute values give only a logarithm. Write $R(y,y')=R(y,y)+[R(y,y')-R(y,y)]$. Then $\sum_{y'}\xi^3(\partial_{\mu'}C\partial_\mu^*)(y'-y)=0$ exactly (the sum over $y'$ of a forward difference in $y'$ of a decaying function), and $|R(y,y')-R(y,y)|\le C'\rho(1+|\log(\xi+\rho)|)$ for $y'\ne y$, by integrating the pointwise bound $|\partial R(y,z)|\le C(1+|\log(\xi+|z-y|)|)$ of \cref{sup4:lem:R} along a monotone path from $y$ to $y'$ (the integral is $\le C(\xi+\rho)(2+|\log(\xi+\rho)|)\le C'\rho(1+|\log(\xi+\rho)|)$ since $\rho\ge\xi$; the cruder bound $|y'-y|_1\sup|\partial R|$ would give a $\log(1/\xi)$). This leaves $\sum\xi^3\rho(1+|\log(\xi+\rho)|)(\xi+\rho)^{-3}e^{-\delta\rho}\le C$. The tadpoles with $R$ are bounded by $|R(y,y)|\le C$ and $\xi|(R\partial^*)(y,y)|\le C\xi(1+|\log\xi|)$. For the first moment there is an extra factor $\rho$, and all four kinds are bounded by absolute values; the fourth gives $\sum\xi^3\rho(\xi+\rho)^{-3}e^{-\delta\rho}\le C$.
\end{proof}

\begin{remark}\label{sup4:rem:S2111}
Ba{\l}aban's ``If at least one propagator $G_{j_0}(0)$ is replaced by $G_{j_0}(0)(1-m^2-aP)C^\xi$, then we get a convergent expression'' \cite[p.~441]{Ba3} is true, but for the term $R\cdot\partial C\partial^*$ of the bubble it needs the cancellation $\sum_{y'}\partial_{\mu'}(\cdots)=0$ above, not power counting with absolute values: in position space $R$ is bounded but does not vanish at coincident points. The cancellation uses only the bounds of \cref{sup4:lem:R} on the $R$-line and the forward-difference structure of the other line.
\end{remark}

\begin{lemma}[{the subtraction identities of the $D=0$ vector and $\phi\phi A$ classes, \cite[(3.31)--(3.38)]{Ba3}}]\label{sup4:lem:S2112}
\textup{(a)} For the remaining divergent two-vector-leg graphs (orders $e_k^4$, $e_k^2\lambda$, with counterterm insertions; at most four vertices by \cref{sup4:lem:S211}), \cite[(3.31)]{Ba3} is an exact identity: the H\"older term has degree $+\alpha$, and the local term has coefficient $\sum_{x'}\eta^d\Pi_{\mu\mu'}(x,x')$. After regrouping and rescaling, the all-$C^\xi$ part of this coefficient vanishes by \cref{sup4:lem:S216}(d) for the full class of that order (the classes ``containing all possible connected graphs \dots\ of a fixed order'' \cite[p.~431]{Ba3}); the localisation functions of internal photon legs are first moved to $x$ by Taylor in $g$, which costs a degree $+1$ (\cite[p.~444]{Ba3}, \cite[(3.36)]{Ba3}). The terms with at least one $R$-line gain at least $1$ in the superficial degree: $|R|\le C$ against $(\xi+\rho)^{-1}$, $|\partial R|\le C(1+|\log|)$ against $(\xi+\rho)^{-2}$, $|\partial R\partial^*|\le C(\xi+\rho)^{-1}$ against $(\xi+\rho)^{-3}$. \textup{(b)} For the $\phi\phi A$ class: \cite[(3.34)]{Ba3} is an exact rewriting (transport the legs with the lowest indices to the vertex with the highest-index leg, through H\"older differences); the localisation functions are moved to the vertex $x'$ by Taylor in $g$, costing degree $+1$ (\cite[(3.36)]{Ba3}); after regrouping, rescaling and $C$-replacement the all-$C$ zero-momentum coefficient $\sum_{y,y''}\xi^{2d}\Gamma''_\mu$ vanishes for the full class of each order by \cref{sup4:lem:S217}(a), of which Ba{\l}aban's class-by-class vanishings are instances or consequences: \cite[(2.21a)]{Ba3} by \cref{sup4:lem:S212}; \cite[(2.18)]{Ba3} ``the same expressions with opposite signs''; \cite[(2.19)]{Ba3} by \cite[(3.37)]{Ba3}; \cite[(2.20)]{Ba3} by \cite[(3.38)]{Ba3}; \cite[(2.21c)]{Ba3} by translation invariance; \cite[(2.21b), (2.21g)]{Ba3} ``with the proper counterterm''. The $R$-terms gain $\ge1$ as in (a).
\end{lemma}

\begin{proof}
The identities are Taylor's formula (\cref{sup4:lem:S213}) at order $0$ and the definitions; the vanishings are \cref{sup4:lem:S216}(d) and \cref{sup4:lem:S217}; the gains of the $R$-terms are \cref{sup4:lem:R} against \cref{sup4:lem:C}.
\end{proof}

\paragraph{What remains of S21, and where it is proved.} The identities above reduce \cite[(3.5)]{Ba3} to three statements about the regrouped totals: that the $R$-terms of the two-loop $D=0$ classes (the superficial gain of \cref{sup4:lem:S2112} made into a bound with the inner subtractions) are bounded uniformly in $\xi$; that the legs can be transported across the subgraphs \cite[(2.4)]{Ba3} so that each is either made convergent by a H\"older factor or admits the summation by parts \cite[(2.8)]{Ba3} without a derivative on an external field (\cite[p.~443]{Ba3}); and the induction of \cite[pp.~444--445]{Ba3} over the orderings, with the vacuum classes cancelled exactly by the vacuum counterterms \cite[(1.24)]{Ba3} up to the image terms of \cref{sup4:lem:F2} (vacuum graphs of at most six vertices, \cref{sup4:lem:S211}). These are the renormalisation theorem of the main paper (Step~3(D$_4$)), which takes \cref{sup4:lem:S211,sup4:lem:S212,sup4:lem:S213,sup4:lem:S214,sup4:lem:S215,sup4:lem:S216,sup4:lem:S217,sup4:lem:S218,sup4:lem:S219,sup4:lem:S2110,sup4:lem:S2111,sup4:lem:S2112} and \cref{sup4:cor:Z1,sup4:lem:C,sup4:lem:K,sup4:lem:R} as its inputs. In it, every inner processing is an exact rewriting (\cref{sup4:lem:S213,sup4:lem:S214}) and the counterterm graphs carry the same lines (\cite[p.~417]{Ba3}), so the total over the regrouped index range is the original total and \cref{sup4:lem:S216}(d), \cref{sup4:lem:S217} and the exact counterterm cancellation apply to it.
